\begin{table}[htbp]
\caption{Uncharacterised mitochondrial ORFs: conservation across \textit{Pleurotus} against expression.}\label{tab:orfcons}
\footnotesize
\begin{tabular}{@{}lrrrc@{}}
\toprule
ORF & Other genomes & Libraries & Reads & In rRNA block \\
\midrule
CM148777.1\_27 & 10/11 & 12/16 & 584 & no \\
CM148777.1\_14 & 10/11 & 11/16 & 524 & no \\
CM148777.1\_44 & 9/11 & 8/16 & 218 & no \\
CM148777.1\_30 & 9/11 & 5/16 & 100 & no \\
CM148777.1\_21 & 9/11 & 4/16 & 60 & yes \\
CM148777.1\_37 & 9/11 & 3/16 & 29 & yes \\
CM148777.1\_43 & 8/11 & 6/16 & 106 & no \\
CM148777.1\_20 & 7/11 & 5/16 & 104 & yes \\
CM148777.1\_28 & 7/11 & 1/16 & 5 & no \\
CM148777.1\_26 & 7/11 & 0/16 & 0 & no \\
CM148777.1\_33 & 6/11 & 11/16 & 376 & no \\
CM148777.1\_51 & 6/11 & 7/16 & 58 & yes \\
CM148777.1\_11 & 6/11 & 3/16 & 30 & no \\
CM148777.1\_50 & 5/11 & 16/16 & 14,070 & yes \\
\bottomrule
\end{tabular}
\begin{tablenotes}[flushleft]\footnotesize\item No Swiss-Prot hit and no informative Pfam domain. An ORF conserved in most genomes and expressed in most libraries is a candidate gene; one present in a single strain is more likely spurious.\end{tablenotes}
\end{table}
